\documentclass[conference]{IEEEtran}
\IEEEoverridecommandlockouts
% The preceding line is only needed to identify funding in the first footnote. If that is unneeded, please comment it out.
\usepackage{cite}
\usepackage{amsmath,amssymb,amsfonts}
\usepackage{algorithmic}
\usepackage{graphicx}
\usepackage{textcomp}
\usepackage{xcolor}
\usepackage{bm}
\usepackage[T1]{fontenc}
\usepackage{lmodern}
\usepackage{tabularx}
\usepackage{url}
\def\BibTeX{{\rm B\kern-.05em{\sc i\kern-.025em b}\kern-.08em
    T\kern-.1667em\lower.7ex\hbox{E}\kern-.125emX}}
\begin{document}

\title
{GenoID: A Declarative RFC 9562 v8 UUID Composition Framework via Genetic-Algorithm-Style Constraint Repair\\}

\author{
\and
\IEEEauthorblockN{Padala Manohar Maharshi}
\IEEEauthorblockA{\textit{Department of Computer Science \& Engineering(AI)} \\
\textit{Vignan\'s Institute Of Information Technology}\\
Visakhapatnam, India \\ pmanoharmaharshi@gmail.com}
\and
\IEEEauthorblockN{Ouku Jasmine}
\IEEEauthorblockA{\textit{Department of Computer Science \& Engineering} \\
\textit{Vignan\'s Institute Of Information Technology}\\
Visakhapatnam, India \\
oukujasmine@gmail.com}
}

\maketitle

\begin{abstract}
Standard UUID generators cannot embed application-defined structure a shard, a tenant, a monotonic counter inside a standards-compliant identifier without paying rejection sampling's exponential cost. RFC 9562 makes this gap explicit: version 8 reserves 122 bits for exactly this purpose but defines no algorithm to fill them, so every production system that wants a self-describing identifier today hand-rolls a one-off, fixed layout. We close this gap with GenoID, a declarative composition framework that produces valid, collision-free v8 UUIDs from arbitrary constrained fields at O(k) cost instead of $\bm{O(64^k)}$ verified across 100M generated identifiers with zero collisions and a full NIST SP 800-22 and dieharder statistical pass. The result is a portable algorithm for the UUID version RFC 9562 defines but does not implement, backed by formal complexity and entropy-preservation proofs and an explicit adversarial security analysis that separates GenoID's architectural contribution from the unchanged, CSPRNG-inherited randomness quality beneath it.
\end{abstract}

\begin{IEEEkeywords}
UUID, RFC 9562, identifier generation, genetic algorithm, constraint repair, CSPRNG, distributed systems, database indexing.
\end{IEEEkeywords}

\section{Introduction}
A UUID that cannot describe itself forces every system that touches it to keep a side table. Distributed systems increasingly rely on 128-bit UUIDs as primary keys and event identifiers precisely because they can be generated independently, without coordination, across nodes but standard UUIDs are opaque. RFC 9562 (May 2024), which obsoletes RFC 4122, formalizes this space and adds three new versions: v6 and v7 solve the B-tree fragmentation problem of purely random v4 IDs by adding a timestamp, and v8 is left explicitly experimental its 122 non-fixed bits carry no defined semantics, and the standard states plainly that "UUIDv8's uniqueness will be implementation specific and MUST NOT be assumed" [1, §5.8]. The widely used uuid JavaScript package ships no v8() generator at all, noting that "the RFC does not define a creation algorithm for them" [17]. Version 8 is a form with 122 blanks and no instructions for filling them in.

Consider a team running a sharded Postgres cluster who wants every primary key to encode its shard number, so a router can locate the owning node without a lookup table. UUIDv4 gives them randomness but no shard field. UUIDv7 gives them a timestamp but still no shard. Rejection sampling draw a candidate UUID, keep it only if the shard bits happen to land in the allowed set of, say, five values, and redraw otherwise needs up to $O(64^k)$ expected draws for k such constrained fields ($6.9\times10^{10}$ for k=6), which is unusable in production. GenoID generates the same key correctly on the first try, every time, in O(k) work.

Three existing paths fall short of this. UUIDv4 is fully random and carries no structure by design. UUIDv7 fixes a single timestamp layout and nothing beyond it. The closest code-level prior art, \texttt{pg\_uuid\_v8}, hides one encrypted timestamp inside a v4-format UUID — valuable, but a single fixed layout with no multi-field composition. None offers a declarative general mechanism for arbitrary structured fields; a full comparison is given in Section~VI.

GenoID resolves this without rejection by borrowing the vocabulary of genetic algorithms, population, crossover, and mutation, and re-purposing it as an engineering mechanism rather than an optimization search. Two independently CSPRNG-populated candidate UUIDs act as parents; a child is assembled by field-boundary crossover, with each declared field inherited whole from one parent. Any field that violates its constraint is then repaired in place by a deterministic nearest-valid-value mapping, in $O(\text{field length})$ time, with no redraw and no retry loop.

\noindent\textbf{Contributions.}

\begin{itemize}
    \item A declarative RFC~9562 v8 layout abstraction (\texttt{V8Layout}/\texttt{V8Field}) for typed, constrained fields, the algorithm the RFC leaves undefined (Section~II).
    
    \item Field-boundary crossover over two independently pooled CSPRNG parents, proved to preserve min-entropy on random-type fields (Section~III-B).
    
    \item Constraint-guided mutation as repair, proved $O(k)$ per UUID against rejection sampling's $O(64^k)$ (Section~III-A), and validated empirically at 0 violations over 1.5M checks (Section~V-A).
    
    \item An adversarial security analysis that discloses a bounded 256-UUID forward-secrecy window and the by-design metadata leakage of structured layouts, rather than resting on NIST PASS alone (Section~IV).
    
    \item Validation across five axes: composition correctness, collision-freedom at 100M scale, NIST SP~800-22 and \texttt{dieharder} statistical batteries, and multi-runtime throughput against the closest prior art (Section~V), with the resulting limitations disclosed openly (Section~VII).
\end{itemize}

\section{Design}

Think of a UUID's 122 free bits as a form with labeled blanks: some blanks must hold a timestamp, one must hold a shard number drawn from a small allowed set, one must only ever increase. GenoID lets a developer declare that form once, then fills every blank correctly on every generation call — no blank is ever left invalid, and no call is ever discarded.

\subsection{Declarative layout}

A \texttt{V8Layout} declares an ordered set of \texttt{V8Field}s, each with a bit offset, a bit length, a type (\texttt{timestamp-ms}, \texttt{counter}, \texttt{shard}, \texttt{node}, or \texttt{random}), and an optional constraint (allowed set, min/max bound, or monotonic). Two shipped layouts illustrate the DSL: \texttt{dbkey} (48-bit timestamp, 8-bit shard constrained to a five-value allowed set, 16-bit monotonic counter, remaining bits random) and \texttt{multitenant} (12-bit tenant enum, 8-bit region enum, remaining bits random). \texttt{completeLayout} fills any undeclared bit range with random-type filler, so every generated value is a syntactically valid RFC~9562 v8 UUID, version nibble 8 and variant bits fixed, the remaining 122 bits carrying the declared composition.

\subsection{Field-boundary crossover}
Generation begins by populating two independent CSPRNG pools, the ``parents,'' in full: \texttt{crypto.getRandomValues} fills the entire pool buffer before any field-specific value is written, so every field position in both parents starts from fresh, mutually independent CSPRNG bytes. A child UUID is then assembled field by field. For each declared field, one bit of an independently drawn \texttt{fieldSelect} value chooses whether the child inherits that field's bits from parent A or parent B. Because every field in both parents is independently CSPRNG-populated, and the select bit is independent of field contents, crossover never mixes a partial ``good'' value with a partial ``bad'' one within a field; it selects one whole, independently random field value from one of two independent sources.

\subsection{Constraint-guided mutation as repair}
After crossover, \texttt{repairConstraints} performs a single pass over the layout's constrained fields, with no retry loop and no redraw. An allowed-constrained field is mapped to its Hamming-nearest member of the allowed set; a min/max-constrained field is clamped; a monotonic field is clamped to \texttt{max(current, last-seen)}, and the new value is recorded for the next call. Every branch returns a value that satisfies its constraint by construction; membership, clamping, and monotonicity are structural guarantees, not probabilistic outcomes. Section~V-A reports the empirical face of this guarantee: 0 violations, not ``near 0.''

\section{Formal Analysis}
Rejection sampling and repair make fundamentally different bets. Rejection bets that a lucky whole-UUID draw arrives before an unlucky run goes on forever; repair guarantees success on every draw by fixing only the fields that need it. The exponential-versus-linear gap below quantifies that bet, and the entropy argument that follows shows the bet costs nothing in randomness quality.

\subsection{Repair complexity: $O(k)$ versus $O(64^k)$}
\textbf{Setup.} Let a layout declare $k$ constrained fields $f_1, \dots, f_k$, each of bit-length $\ell_i$. Rejection sampling draws a full candidate UUID, checks all $k$ constraints, and, on any failure, discards the entire UUID and redraws from scratch.

Rejection sampling is exponential in expectation. For an allowed-constrained field of length $\ell$ with $a$ valid values out of $2^\ell$, a uniform draw succeeds with probability $a / 2^\ell$. For $k$ independent fields, whole-UUID acceptance probability is the product $\prod (a_i / 2^{\ell_i})$. In the worst case ($a_i = 1, \ell_i = 8$, one allowed value per byte-granular field), this is $2^{-8k}$, giving an expected $2^{8k} = 64^k$ whole-UUID draws, roughly $6.9 \times 10^{10}$ for $k=6$. This expectation is a geometric random variable: a single unlucky run has no upper bound.

\texttt{repairConstraints} is linear in $k$. The repair function iterates the layout's fields exactly once. Each allowed-constrained field costs $O(|\text{allowed}_i| \cdot \ell_i)$, a bounded Hamming-distance search; min/max and monotonic fields cost $O(1)$; writing a repaired field back costs $O(\ell_i)$. Because $|\text{allowed}_i|$ and $\ell_i$ are layout constants independent of $k$, the total cost $\sum O(|\text{allowed}_i| \cdot \ell_i)$ is $O(k)$, linear, with an empirically measured constant factor bounded by 8. Critically, this bound holds unconditionally: the repair function never redraws and never loops more than once over the fields, eliminating rejection sampling's unbounded tail risk entirely.

\textbf{Correctness.} Each repair branch returns a constraint-satisfying value by construction; the proof, not an independent empirical coincidence, is why Section V-A measures exactly 0 mismatches and 0 constraint violations over 1.5M structured-field checks.

\subsection{Entropy preservation under field-boundary crossover}

Shuffling two decks of independently drawn random bits, field by field, does not make the result less random than either deck alone; the claim below states this precisely.

\textbf{Claim.} For a random-type field, field-boundary crossover does not reduce min-entropy relative to drawing that field's bits directly from the CSPRNG, conditioned on the field-select bit being independent of field contents.

\textbf{Argument.} Let $X_A, X_B$ be a field's value in parent A and B, both uniform on $\{0,1\}^\ell$ and mutually independent, since both parents are independently CSPRNG-populated before any field write. Let $S$ be the independent select bit; the child's value is $X_A$ if $S=1$, else $X_B$. For any fixed $v \in \{0,1\}^\ell$:
\[
\Pr[\text{child} = v] = \Pr[S{=}1]\,2^{-\ell} + \Pr[S{=}0]\,2^{-\ell} = 2^{-\ell}.
\]
The child field is itself uniform on $\{0,1\}^\ell$ regardless of $S$'s distribution; crossover is measure-preserving on random-type fields. This matches the empirical NIST SP 800-22 15/15 PASS result on structured layouts' random payload; bias would surface as a monobit or frequency test failure. It also explains why crossover provably cannot improve entropy: it is a no-op on an already-uniform source. For structured (non-random) field types, the claim does not apply and is not claimed; those fields are deterministic or constrained by design and correctly carry zero min-entropy in the accounting of Section~IV-B regardless of crossover.

\textbf{Repair does not increase entropy.} \texttt{repairConstraints} is a deterministic function of its input; by the data-processing inequality, $H(g(X)) \leq H(X)$ for any deterministic $g$, so repair can only hold or reduce entropy on a constrained field. This is the formal grounding for the paper's central honesty claim, restated once more because it matters: GA's contribution in GenoID is architectural, composition and constraint satisfaction, never a randomness enhancer.

\section{Security Analysis}

Statistical randomness alone does not establish security. A NIST SP 800-22 PASS shows the output resembles random noise; it does not show the output is unpredictable to an adversary who can act. This section gives that second, adversarial argument, which the composition mechanism of Section II does not by itself provide.

\subsection{Threat model}
\begin{itemize}
    \item \textbf{Passive observer.} Capability: collects $\leq N$ emitted UUIDs, no process access. Goal: distinguish GenoID from random; infer structure; predict future values.
    
    \item \textbf{Distinguishing attack.} Capability: passive, plus an oracle. Goal: decide whether a given UUID is GenoID- or CSPRNG-derived.
    
    \item \textbf{State compromise.} Capability: reads process memory (in-process CSPRNG pool). Goal: predict future UUIDs until the next pool refill.
    
    \item \textbf{Backward secrecy.} Capability: reads current pool state. Goal: recover past UUIDs.
    
    \item \textbf{Structure inference.} Capability: holds one or more structured UUIDs. Goal: read timestamp/shard/tenant/counter (by design, for routing).
\end{itemize}

The OS CSPRNG (\texttt{crypto.getRandomValues}) is assumed reliable, matching RFC~9562 Section~8.1; every GenoID pool refill is seeded from it.

\subsection{Entropy accounting}

\begin{center}
\footnotesize

\textbf{TABLE I.} Min-Entropy by Generator

\vspace{0.4em}

\begin{tabular*}{\columnwidth}{@{\extracolsep{\fill}}lccc}
\hline
\textbf{Generator} & \textbf{Bits} & \textbf{Secret} & \textbf{Entropy} \\
\hline
UUIDv4 & 128 & 122 random & 122 \\
UUIDv7 & 128 & 74 random & 74 \\
GenoID v8 & 128 & 122 CSPRNG & 122 \\
GenoID dbkey & 128 & 50 random & 50 \\
ULID-v8 & 128 & 74 random & 74 \\
pg\_uuid\_v8 & 128 & Up to 122 random & Up to 122 \\
Math.random & 128 & PRNG state & 0 \\
\hline
\end{tabular*}
\end{center}

\indent UUIDv4 uses \texttt{crypto.randomUUID()} and \texttt{pg\_uuid\_v8} embeds an AES-encrypted timestamp. For the \texttt{dbkey} layout,
$128 - 48\,(\texttt{ts}) - 8\,(\texttt{shard}) - 16\,(\texttt{counter}) - 6\,(\texttt{fixed}) = 50$
so 50 random bits are the only entropy an adversary cannot strip.
\subsection{Adversarial findings}
Non-structured bits are CSPRNG-drawn and computationally indistinguishable from uniform random, matching RFC~9562 Section~8.1 and consistent with the 15/15 NIST PASS of Section~V. Structured layouts leak metadata by design; that is what makes them self-describing and routable. They are explicitly not a confidentiality mechanism. \texttt{pg\_uuid\_v8} is strongest on this particular axis: its AES-encrypted timestamp makes the full 122 bits look random even to a passive observer, a genuinely elegant solution to a problem GenoID does not attempt to solve.

GenoID's in-process pool holds 256 entries. On refill, \texttt{crypto.getRandomValues()} repopulates all 256 independent of prior state. An adversary who reads the pool buffer can therefore predict at most 256 future UUIDs, after which the stream is unpredictable again, a bounded, documented forward-secrecy window, not a total break. Unpooled generators such as UUIDv4 have strictly better forward secrecy, since there is no pool buffer to steal, at the cost of the throughput that pooling exists to recover; this is a deliberate trade-off, disclosed rather than hidden.

Past pool buffers are overwritten and not retained, so reading the current pool state does not recover past UUIDs. GenoID provides backward secrecy. \texttt{Math.random()} (Xorshift128+) does not: its full generator state is recoverable from a handful of consecutive outputs.
\subsection{Honest security-class summary}
\begin{center}
\scriptsize
\textbf{TABLE II.} Security Class by Generator

\vspace{0.3em}

\begin{tabular*}{\columnwidth}{@{\extracolsep{\fill}}lccc}
\hline
\textbf{Generator} & \textbf{Class} & \textbf{Caveat} \\
\hline
v4 & High & §8.1 security; no pool \\
v7 & High* & Ts leak (§8.2) \\
GenoID v8 & High & $\leq$256-UUID FS window \\
GenoID-struct. & High (comp.) & 50-bit entropy; ts/shard/ctr leak \\
\texttt{pg\_uuid\_v8} & High & AES ts, indistinguishable \\
\texttt{Math.random()} & Insecure & Reversible state, 0-bit \\
\hline
\end{tabular*}
\end{center}

\indent No formal cryptographic reduction proof is claimed here. This analysis is entropy accounting plus adversarial reasoning, not a reduction to a hardness assumption. Such a proof would require modeling the OS CSPRNG as a random oracle, which is outside the scope of this paper and is disclosed as such rather than glossed over.
\section{Evaluation}
Every number below is reproducible from a single command, \texttt{bun run bench}, and is also produced as a CI artifact across a nine-job matrix spanning Ubuntu~24.04, macOS~14, and Windows~2025, each on Bun (latest), Node~22 LTS, and Deno~2.9.3, covering seven distinct runtime$\times$OS configurations.

\subsection{Composition correctness}
Over 1.5M structured-field checks, GenoID reports 0 mismatches and 0 constraint violations, the empirical counterpart to the by-construction correctness proof of Section~III-A.

\subsection{Repair vs.\ rejection cost}
Measured GA repairs per UUID scale as $\approx k$, confirming the $O(k\cdot8)$ bound of Section~III-A against rejection sampling's $O(64^k)$.

\subsection{Collision and uniformity}
Across 2M generated UUIDs, every tested generator, including \texttt{genoid-structured}, reports 0 collisions, with a maximum uniformity deviation of 0.0053 on the random payload (uniformity is measured on the random payload only, not the whole UUID, since structured layouts have a constant leading timestamp byte). At full scale, 100M UUIDs, via an open-addressing 128-bit hash set fanned across all cores, every generator again reports 0 observed collisions, against a theoretical 50\%-collision birthday bound of roughly $2.7\times10^{18}$. This scale is a sanity check against implementation bugs, not a claim to have exhausted the collision space. No UUID paper tests to $2.7\times10^{18}$, and this one does not pretend to.

\subsection{NIST SP 800-22 and dieharder}
All 15 NIST SP~800-22 sub-tests pass on three structured layouts (\texttt{struct-dbkey}, \texttt{multitenant}, and \texttt{eventsourcing}). An extended \texttt{dieharder} battery, including \texttt{birthdays}, \texttt{rank\_32x32}, \texttt{dna}, \texttt{count\_1s\_str}, \texttt{parking\_lot}, \texttt{runs}, \texttt{sts\_monobit}, and \texttt{sts\_serial}, comprising 38 sub-tests, 5 independent trials each, and 100M-bit samples, reports 152/152 PASS (modal assessment per sub-test) across four generator variants: native UUIDv4, \texttt{rawv8} (RFC~9562 v8 without GA), \texttt{genoid-v8} (GA + pooling), and \texttt{struct-dbkey} (structured). The GA and non-GA variants pass identically, which is the empirical counterpart to the Section~III-B proof: GA neither degrades nor improves statistical quality on random-type fields.

\subsection{Throughput and baseline comparison}
\begin{center}
\scriptsize
\textbf{TABLE III.} Throughput (ops/sec, mean of 10 trials, 95\% CI within $\pm$5\%), Ubuntu/Bun column shown; full seven-cell matrix in the repository.

\vspace{0.3em}

\begin{tabular}{lll}
\hline
\textbf{Generator} & \textbf{Ops/sec (Ubuntu, Bun)} & \textbf{NIST} \\
\hline
v4-native & 11.85M & -- \\
v7-custom & 5.74M & -- \\
genoid-v8 (pooled) & 10.12M & -- \\
pg\_uuid\_v8 & 0.85M & 15/15 \\
ulid-v8 & 0.93M & 15/15 \\
snowflake & 3.01M & -- \\
genoid-structured (dbkey) & 0.67M & 15/15 \\
\hline
\end{tabular}
\end{center}

Head-to-head against the closest code-level prior art, \texttt{pg\_uuid\_v8} ($n=2$M): both report 0 collisions; \texttt{genoid-structured}'s uniformity deviation is 0.0051 against \texttt{pg\_uuid\_v8}'s 0.0066; \texttt{pg\_uuid\_v8} runs about $1.7\times$ faster, since cheap XOR steganography beats GA-based repair on raw speed, but it is fixed-layout, timestamp only, while GenoID is declarative and supports arbitrary multi-field composition. Both pass NIST. GenoID trades some throughput for composition flexibility; it does not trade away statistical quality to get it.

A secondary finding concerns the runtime, not the algorithm. Node's \texttt{crypto.getRandomValues()} carries markedly higher per-call overhead than Bun's or Deno's, so generators calling it once per UUID, \texttt{v7}, \texttt{ULID}, \texttt{pg\_uuid\_v8}, \texttt{ULID-v8}, and \texttt{KSUID}, run $3$--$13\times$ slower on Node than on Bun or Deno on comparable OSes. GenoID's pooled design, at 0.0039 CSPRNG calls per UUID, stays within roughly $1.5\times$ across runtimes, a property of amortized pooling, orthogonal to the composition contribution this paper makes.

\subsection{Additional deployability checks}
Concurrent generation via \texttt{worker\_threads} fan-out shows 0 cross-worker collisions and 0 constraint violations. A 100k-ID SQLite B-tree insertion test shows every tested ID type produces a clean B-tree, zero freelist growth, with sortable IDs matching or exceeding random-insert throughput; this is an index-locality proxy, not a claim about production database engines (Section~VII returns to this limitation directly). A Playwright-driven cross-engine browser test, Chromium, Firefox, and WebKit, reports zero browser errors and 0 collisions on all three engines.

\section{Related Work}
\subsection{UUID standards: RFC 4122 to RFC 9562}
UUIDs are 128-bit identifiers, originally standardized as RFC~4122 (2005, versions~1--5), obsoleted by RFC~9562 (IETF Standards Track, May~2024), which adds versions~6, 7, and~8~\cite{1}. UUIDv4 carries 122 bits of CSPRNG randomness with fixed version/variant bits, simple and unpredictable, but with poor B-tree locality due to its lack of order. UUIDv7 prepends a 48-bit big-endian Unix-millisecond timestamp, giving 74 random bits and solving v4's index-fragmentation problem while remaining drop-in compatible with existing v4 storage~\cite[\S5.7]{1}. UUIDv8 is explicitly experimental: only the version and variant bits are defined, and the RFC supplies no creation algorithm for the remaining 122~\cite[\S5.8]{1}.
\subsection{Sortable and structured identifiers}
RFC~9562's own motivation section surveyed sixteen prior implementations while drafting the standard, ULID, LexicalUUID, Snowflake, Flake, ShardingID, KSUID, Elasticflake, FlakeID, Sonyflake, orderedUuid, COMBGUID, SID, pushID, XID, ObjectID, and CUID~\cite{1}. Representative members: ULID~\cite{2} (48-bit timestamp + 80-bit randomness, Base32, lexicographically sortable); KSUID~\cite{3} (32-bit second-resolution timestamp + 128-bit randomness, tuned for distributed log ordering); Snowflake and derivatives~\cite{4} (timestamp $|$ worker $|$ sequence, k-sortable but coordination- and clock-dependent); TypeID~\cite{5} (typed prefix over a UUIDv7 payload); and COMBGUID (hybrid UUID with an embedded timestamp for index locality). Empirically, time-ordered identifiers measurably improve database behavior versus random v4, reported gains include roughly 35\% faster inserts and 22\% smaller indexes at 10M rows, and Shopify reporting that switching to ULID halved \texttt{INSERT} duration on a high-throughput MySQL table~\cite{15,16}. Recent academic work directly compares UUIDv4/v7/ULID for distributed systems on collision probability, network overhead, and generation speed~\cite{12}, and a further study benchmarks ULID against UUIDv4/v7 atop Kafka and PostgreSQL~\cite{13}. This is a well-explored field, its gap is not effort but scope: it has explored time-ordered composition thoroughly, almost exclusively via fixed, hand-coded layouts, and none of it offers a declarative framework for arbitrary multi-field v8 composition.
\subsection{Steganographic UUIDs — closest prior art}
\texttt{pg\_uuid\_v8}~\cite{8}, a PostgreSQL C extension (PGXN, May~2026), is the closest prior art, and a clever one. It generates v4-format-compliant UUIDs while embedding an encrypted microsecond timestamp, XOR, AES-128, or AES-256, in the random portion, recoverable via a dedicated extraction function and indexable through a functional index, resolving the same performance/privacy trade-off as v7 through steganography rather than a plaintext timestamp. GenoID shares \texttt{pg\_uuid\_v8}'s core instinct, put useful structure inside a standards-compliant UUID, but differs in scope: \texttt{pg\_uuid\_v8} offers one fixed stego layout with no multi-field abstraction and no repair mechanism, while GenoID is a declarative, multi-field DSL packaged as a portable library (browser, Node, and CI-verified) rather than a database-specific extension. The two are complementary, not competing: a system wanting encrypted-timestamp indexability inside Postgres should still reach for \texttt{pg\_uuid\_v8}.


\begin{center}
\scriptsize
\textbf{TABLE IV.} Capability comparison with the closest prior art.

\vspace{0.3em}

\begin{tabularx}{\columnwidth}{@{}lXX@{}}
\hline
\textbf{Capability} & \textbf{pg\_uuid\_v8} & \textbf{GenoID} \\
\hline
UUID form &
v4-compatible steganographic &
RFC~9562 v8 (custom) \\

Layout &
Single fixed stego layout &
Declarative \texttt{V8Layout/}\allowbreak\texttt{V8Field} DSL \\

Multi-field support &
No &
Yes (timestamp, shard, counter, tenant, \ldots) \\

Repair mechanism &
Not applicable &
Constraint-guided mutation as repair \\

Platform &
PostgreSQL extension, code-only &
Portable TS library, browser + Node + CI \\
\hline
\end{tabularx}
\end{center}

\subsection{Genetic algorithms and evolutionary computation}
A genetic algorithm (GA) is a population-based optimizer applying selection, crossover, and mutation, inspired by natural selection, and surveyed comprehensively as a general optimization family~\cite{9,10,11}. GAs are applied broadly across engineering and, recently, in synergy with large language models~\cite{19}. A disclosed literature search across Semantic Scholar, arXiv, OpenAlex, and general web sources, re-verified in a July~2026 adversarial recheck aimed specifically at falsifying this claim, surfaces no academic paper applying genetic or evolutionary algorithms to UUID or identifier generation. Targeted queries combining GA terminology with identifier-generation terminology return two disjoint literatures with no intersection found; a parallel patent search finds GA-machine and genetic-programming patents (e.g., US5343554, US5970487, US6360191B1) alongside three unrelated identifier-generation patents that use hashing, counters, or coordination protocols, none combines the two. GenoID is accordingly not a new GA technique, but the first application of GA-style operators to the composition of RFC~9562 version~8 UUID payloads, and Section~III-B's proof is exactly the honest accounting this claim demands. GA's value here is architectural, not an entropy improvement, and the paper says so rather than implying otherwise.
\subsection{High-throughput secure generation}
v4/v8 generation throughput is bounded by the cost of drawing from the OS CSPRNG. Prior work recovers throughput by reusing a single CSPRNG source and amortizing draws. Standard Java guidance favors reusing one \texttt{SecureRandom} instance over per-call construction, and a pooled/batch Go UUID library reports roughly a $14\times$ speedup for pooled versus stateless v4 generation~\cite{14}. PostgreSQL~18 now ships a native \texttt{uuidv7()}~\cite{18}. GenoID adopts the same pooling principle at the algorithm level. One CSPRNG call fills a pool covering 64 UUIDs' worth of bytes, subsequently composed via byte-level crossover and mutation, combining pooling throughput with v8 composition flexibility.

\section{Threats to Validity}
We disclose these openly, ranked by how strongly we expect a reviewer to press on each. An unacknowledged limitation invites a hostile review; an acknowledged one invites a collaborative one.

\textbf{External validity,} the strongest limitation. This is a single-language, single-ecosystem (TypeScript) implementation. The composition algorithm is language-agnostic, byte-array bit operations translate directly, but no port to Go, Rust, or Java exists yet to confirm the complexity and throughput claims transfer. Benchmarks run on GitHub Actions-hosted CI runners (shared vCPUs, no NUMA, capped burst credit), so absolute throughput numbers should be read as relative comparisons across generators on the same runner, not production capacity-planning figures. The SQLite B-tree test uses a single engine and a uniform key distribution; real production skew, hot shards, bursty tenants, is not modeled and is left as future work.

\textbf{Construct validity.} ``Statistical quality'' (NIST PASS/FAIL) is scoped explicitly as necessary but not sufficient for security, which is exactly why Section~IV exists as an independent adversarial argument rather than a restatement of the NIST result. The novelty construct, ``no prior GA-for-identifier work,'' rests on a disclosed, re-verified, but non-exhaustive search; absence of evidence is not evidence of absence, though the breadth of the search (multi-database, patent, and web, with the closest adjacent papers explicitly cited and differentiated) mitigates this substantially rather than eliminating it.

\textbf{Internal validity.} Benchmark confounds, JIT warm-up, GC pauses, are mitigated by reporting mean $\pm$ standard deviation over 10 trials with 95\% confidence intervals and a Welch t-test with Cohen's $d$, rather than single-shot numbers. Two implementation bugs (32-bit truncation; single-parent population) were caught by the test suite before any result was reported. NIST's full 15-test battery runs unconditionally on every sample; the \texttt{dieharder} subset is an explicitly disclosed curation for CI time budget, not a hidden one.

\textbf{Conclusion validity.} Throughput comparisons use Welch t-tests and Cohen's $d$ rather than bare point estimates. NIST's per-test $\alpha=0.01$ is followed without a stricter family-wise (Bonferroni-style) correction across the 15-test battery, which a security venue may reasonably prefer. ``0 collisions'' is reported alongside the theoretical birthday-bound expectation, not treated as ``0 probability.''

None of these four items is hidden; each is the honest boundary of a claim made earlier in the paper, not a surprise saved for a rebuttal.

\section{Applications}
\begin{itemize}
    \item \textbf{Sharded databases.} Embed the shard ID directly in the primary key; a router locates the owning node without a lookup table.

    \item \textbf{Multi-tenancy.} Carry a tenant ID in the identifier for prefix-based isolation and row-level security enforcement.

    \item \textbf{Event sourcing.} A monotonic counter combined with a timestamp field yields globally ordered, collision-free event identifiers.

    \item \textbf{Sortable time-series data.} Timestamp bits provide chronological order while remaining fields carry composable application data.

    \item \textbf{Operational debuggability.} Declared fields are readable directly from the hex representation, letting operators inspect shard, tenant, or sequence without decoding logic.
\end{itemize}
\section{Conclusion and Future Work}
RFC~9562 leaves version~8's 122 bits undefined on purpose; GenoID fills them with a declarative, portable algorithm instead of a one-off hand-rolled layout. Its central move, genetic-algorithm-style crossover and mutation repurposed as $O(k)$ constraint repair rather than a randomness search, carries a formal complexity proof, a formal entropy-preservation proof, and broad empirical validation: composition correctness, collision-freedom at 100M scale, NIST and \texttt{dieharder} statistical batteries, and throughput benchmarking against the closest prior art. An explicit adversarial security analysis exists specifically so ``passes NIST'' is never mistaken for ``is secure,'' and it discloses a bounded 256-UUID forward-secrecy window and the by-design metadata leakage of structured layouts rather than eliding them.

The paper's own threats-to-validity analysis names the highest-value next step plainly: confirming that the $O(k)$ complexity and throughput characteristics transfer to a non-JavaScript implementation, Go, Rust, or Java, is the single most consequential piece of future work, alongside a formal cryptographic reduction proof for the pooled-CSPRNG construction and an evaluation of B-tree behavior under realistic, skewed production key distributions rather than uniform synthetic ones.

\begin{thebibliography}{00}

\bibitem{1}
K. Davis, B. Peabody, and P. Leach, ``Universally Unique IDentifiers (UUIDs),'' RFC 9562, Internet Engineering Task Force (IETF), May 2024. [Online]. Available: \url{https://www.rfc-editor.org/rfc/rfc9562}

\bibitem{2}
A. B. Peabody, ``ULID Specification.'' [Online]. Available: \url{https://github.com/ulid/spec}

\bibitem{3}
Segment, ``KSUID.'' [Online]. Available: \url{https://github.com/segmentio/ksuid}

\bibitem{4}
Twitter, ``Snowflake.'' [Online]. Available: \url{https://github.com/twitter-archive/snowflake}

\bibitem{5}
Jetify, ``TypeID.'' [Online]. Available: \url{https://github.com/jetify-com/typeid}

\bibitem{6}
rs, ``xid.'' [Online]. Available: \url{https://github.com/rs/xid}

\bibitem{7}
J. Nilsson, ``The cost of GUIDs as primary keys,'' InformIT, 2002.

\bibitem{8}
ineron, ``pg\_uuid\_v8: PostgreSQL extension for steganographic UUIDs,'' Version 1.0.0, May 2026. [Online]. Available: \url{https://github.com/ineron/pg_uuid_v8}

\bibitem{9}
B. Alhijawi and A. Awajan, ``Genetic algorithms: Theory, genetic operators, solutions, and applications,'' \emph{Evolutionary Intelligence}, vol. 17, pp. 1245--1256, 2024.

\bibitem{10}
T. Bäck, \emph{Evolutionary Algorithms in Theory and Practice}. New York, NY, USA: Oxford University Press, 1996.

\bibitem{11}
D. B. Fogel, \emph{Evolutionary Computation: Towards a New Philosophy of Machine Intelligence}. Piscataway, NJ, USA: IEEE Press, 1995.

\bibitem{12}
``A comparative analysis of identifier schemes: UUIDv4, UUIDv7, and ULID for distributed systems,'' arXiv:2509.08969, Sep. 2025.

\bibitem{13}
nimakarimiank, ``uids-comparison.'' [Online]. Available: \url{https://github.com/nimakarimiank/uids-comparison}

\bibitem{14}
pscheid92, ``uuid: Go implementation of RFC 9562 with V8 and pooled/batch APIs.'' [Online]. Available: \url{https://pkg.go.dev/github.com/pscheid92/uuid}

\bibitem{15}
Rich Dev Tools, ``UUID v4 vs v7 vs ULID: What the public benchmarks actually show,'' 2026.

\bibitem{16}
remove.sh, ``UUID v4 vs v7 vs ULID: Which should you use in 2026?,'' 2026.

\bibitem{17}
uuidjs, ``uuid.'' [Online]. Available: \url{https://github.com/uuidjs/uuid}

\bibitem{18}
PostgreSQL Global Development Group, ``UUID functions,'' PostgreSQL 18 Documentation. [Online]. Available: \url{https://www.postgresql.org/docs/current/functions-uuid.html}

\bibitem{19}
``Evolutionary Computation and Large Language Models: A Survey of Methods, Synergies, and Applications,'' arXiv:2505.15741, May 2025.

\end{thebibliography}

\end{document}
